\section{Conclusion and Outlook}
\quad In this work, we developed a computable linear-response framework for noisy network dynamics near stable fixed points. Linearizing the stochastic dynamics around the noise-free fixed point reduces the problem to a Jacobian matrix ${\bf Q}$ and a steady-state covariance matrix ${\bf K}_0$ satisfying a Lyapunov equation. Within this structure, we derived first-order response formulas for perturbations of intrinsic node parameters, connection weights, and noise strengths. The main result is an operational theory for the equal-time covariance response that can be evaluated directly from the unperturbed linearized network and applied to both equilibrium and nonequilibrium steady states with additive Gaussian white noise.\par
The results show that linear response remains predictive for fluctuation observables even when detailed balance is broken. In particular, the covariance response gives a direct way to quantify how local perturbations spread through network-wide fluctuations and correlations. Our Langevin simulations confirmed the steady-state predictions in undirected equilibrium networks and in nonequilibrium networks generated by asymmetric couplings or heterogeneous noise strengths. Tests of multiple simultaneous perturbations further showed that the first-order responses remain additive within the linear regime. The removal examples showed where the first-order approximation begins to break down and thus point to the need for a nonlinear response theory for large perturbations and node removal. Since local perturbations propagate through correlations across the whole network, it will also be important to clarify more systematically how the response depends on topology.\par
We also extended the framework to time-dependent perturbations and obtained explicit time-dependent response formulas for the mean state, deterministic force, and equal-time covariance. This extension shows that the same formalism can describe not only the difference between two steady states, but also the full relaxation process induced by a driving protocol. In this picture, a parameter perturbation drives the system from one steady state to a new one, and the transient response describes how the system moves between them. This viewpoint suggests a natural connection to nonequilibrium steady-state thermodynamics and to the Hatano-Sasa description of transitions between steady states \cite{10.1143/PTPS.130.29,PhysRevLett.86.3463,doi:10.1073/pnas.0406405101}. Another natural extension is to relax the assumption of additive white noise with diagonal covariance. In many realistic systems, fluctuations are temporally correlated or shared across components, leading to colored noise \cite{PhysRevLett.84.3017, PhysRevResearch.5.033208, qrvl-mng6} and noise matrices with off-diagonal elements. Extending the framework to such cases would broaden its applicability and clarify how temporal memory and cross-component noise correlations reshape the response. For nonequilibrium steady states, another direction is to include the energetics of stochastic trajectories \cite{10.1143/PTPS.130.17} and connect the present covariance-based formalism to excess dissipation, housekeeping dissipation, and entropy production \cite{Seifert_2010, Seifert_2012, Mehl_2010}. Since the present theory already gives explicit conjugate-force and time-dependent response kernels in terms of ${\bf Q}$, ${\bf K}_0$, and the perturbation protocol, it should be possible to compute the linear response of these energetic observables directly for noisy networks. Finally, because the current theory is built around fluctuations near stable fixed points, an important next step is to generalize it to oscillatory regimes such as stable limit cycles \cite{PhysRevE.108.064114}.
